
\addcontentsline{toc}{section}{Introduction}
\section*{Introduction}

Ce chapitre illustre la mise en œuvre pratique de notre technique de sécurisation au moyen d'une expérimentation exhaustive avec le modèle ResNet50, dans le cadre d'un système de test pour la détection d'objets dangereux pour les forces de l'ordre. Le but est de prouver la viabilité et l'efficience de l'architecture en cinq zones proposée  au Chapitre 2. 



\section{Présentation du Cas expérimental et choix du modèle.}

\subsection{Objectif de l'expérience}

L'objectif principal ici est de comparer deux approches de développement de modèles d'intelligence artificielle :

\begin{enumerate}
    \item \textbf{Approche Baseline (vulnérable)} : Un modèle ResNet50 entraîné de manière standard via transfer learning, représentant les méthodes usuelles de développement IA sans considération particulière pour la sécurité adversariale.

    \item \textbf{Approche Secured (robuste)} : Un modèle ResNet50 identique architecturalement, mais entraîné avec des défenses multicouches intégrant adversarial training selon la méthode TRADES, differential privacy, et des mécanismes de détection d'anomalies.
\end{enumerate}


\begin{figure}[H]
    \centering
    \includegraphics[width=1\linewidth]{images/Architecture POC.png}
    \caption{Architecture du Cas expérimental : comparaison entre le modèle Baseline et le modèle Secured.}
    \label{fig:enter-label}
\end{figure}


    Cette comparaison permettra de quantifier précisément :
\begin{itemize}
    \item L'\textbf{amélioration de robustesse} apportée par les mécanismes de défense proposés
    \item Le \textbf{trade-off} entre performance sur données propres (clean accuracy) et robustesse adversariale
    \item La \textbf{faisabilité pratique} de déploiement de modèles robustes en production

\end{itemize}

\paragraph{}

Notre expérimentation vise à atteindre de façon spécifique les objectifs suivants :

\begin{enumerate}
    \item Implémenter l'architecture de sécurisation en 5 zones sur un cas réel
    \item Comparer quantitativement un modèle ResNet50 standard versus sécurisé
    \item Démontrer l'amélioration de la robustesse face aux attaques adversariales
    \item Valider la faisabilité technique
    \item Produire un prototype fonctionnel déployable
\end{enumerate}


\subsection{Scénario d'application : Détection d'objets dangereux}

Le scénario choisi pour cette expérience est la \textbf{détection binaire d'objets dangereux versus objets sûrs}, un cas d'usage critique dans les systèmes de maintien de l'ordre où les erreurs de classification peuvent avoir des conséquences graves.

\subsubsection{Contexte opérationnel}

Dans les systèmes de sécurité modernes, les algorithmes de vision par ordinateur sont déployés pour identifier automatiquement des objets dangereux dans des flux vidéo de surveillance, des contrôles d'accès, ou lors de fouilles automatisées. La compromission de ces systèmes par des attaques adversariales pourrait permettre à un adversaire de faire passer un objet dangereux pour un objet sûr, créant un risque de sécurité majeur.

\subsubsection{Classification binaire}

Le problème est formulé comme une tâche de \textbf{classification binaire} :
\begin{itemize}
    \item \textbf{Classe 0 (safe)} : Objets sûrs ne présentant pas de menace
    \item \textbf{Classe 1 (dangerous)} : Objets dangereux nécessitant une intervention
\end{itemize}

Cette simplification binaire, bien que réductrice par rapport à un système de production réel multi-classes, permet de :
\begin{itemize}
    \item Concentrer l'analyse sur la \textbf{robustesse adversariale} plutôt que sur la complexité taxonomique
    \item Faciliter l'\textbf{interprétation des résultats} de sécurité (attack success rate, robustness degradation)
    \item Valider la \textbf{méthodologie de sécurisation} dans un cadre contrôlé avant généralisation
\end{itemize}


\subsection{Justification du choix de ResNet50}

Le choix de l'architecture \textbf{ResNet50} \cite{he2016deep} pour cette PoC repose sur plusieurs critères scientifiques et pratiques.

\subsubsection{Architecture et capacité}

ResNet50 est un réseau de neurones convolutif profond de 50 couches basé sur le mécanisme de \textbf{connexions résiduelles} (residual connections) introduit par He et al. (2016) \cite{he2016deep}.



Cette architecture présente les caractéristiques suivantes :
\begin{itemize}
    \item \textbf{Profondeur} : 50 couches convolutives permettant l'extraction de features hiérarchiques complexes
    \item \textbf{Paramètres} : Approximativement 25 millions de paramètres (25.6M exactement)
    \item \textbf{Connexions résiduelles} : Permettent de résoudre le problème de gradient vanishing dans les réseaux très profonds

\end{itemize}



\subsubsection{Vulnérabilité adversariale documentée}

La littérature scientifique a largement documenté la \textbf{vulnérabilité de ResNet50 aux attaques adversariales} :

\begin{table}[h]
\centering
\caption{Robustesse adversariale de ResNet50 standard (littérature)}
\label{tab:resnet50_vulnerability}
\begin{tabular}{|l|c|c|}
\hline
\textbf{Attaque} & \textbf{Clean Accuracy} & \textbf{Adversarial Accuracy} \\ \hline
Aucune & 76.15\% & 76.15\% \\ \hline
FGSM ($\varepsilon=0.031$) & 76.15\% & $\sim$10-20\% \\ \hline
PGD ($\varepsilon=0.031$, 10 iter) & 76.15\% & $<$10\% \\ \hline
\textbf{Robustness Gap} & -- & \textbf{$>$65\%} \\ \hline
\end{tabular}
\end{table}

Cette vulnérabilité connue en fait un excellent candidat pour démontrer l'efficacité des méthodes de sécurisation proposées dans ce mémoire.


\subsubsection{Adaptabilité au scénario}

ResNet50 s'adapte parfaitement à notre scénario de détection binaire :

\begin{itemize}
    \item \textbf{Transfer Learning} : La dernière couche fully-connected peut être remplacée pour passer de 1000 classes ImageNet à 2 classes (safe/dangerous), tout en conservant les features génériques apprises sur ImageNet
    \item \textbf{Fine-tuning flexible} : Les couches convolutives pré-entraînées peuvent être gelées (pour transfer learning rapide) ou fine-tunées (pour adaptation au domaine spécifique)
    \item \textbf{Taille raisonnable} : $\sim$280MB en format PyTorch (.pth), permettant un déploiement sur des systèmes embarqués ou des serveurs sans GPU haute performance
    \item \textbf{Inference temps réel} : Capable de traiter des images 224×224 pixels à plus de 30 FPS sur GPU moderne, et en temps quasi-réel sur CPU
\end{itemize}




